\section{The logarithmic variational problem}\label{sec:log}

Set
\[
 t=\log r,
\]
and let $\bar\nu$ be the pushforward of $\nu$ under $r\mapsto t$.  Introduce the tilted positive measure $\sigma$ on $\R$ by
\begin{equation}\label{eq:sigma-def}
 \dd\sigma(t):=\ee^t\,\dd\bar\nu(t).
\end{equation}
Then
\[
 A[\sigma]:=\int_\R \ee^{-t}\,\dd\sigma(t)=1,
 \qquad
 B[\sigma]:=\int_\R \ee^t\,\dd\sigma(t)=\int_0^\infty r^2\,\dd\nu(r).
\]

For $t\ge u$, direct substitution gives
\[
 K_3(\ee^t,\ee^u)\,\ee^{-t}\ee^{-u}
 =\frac12\left(\ee^{t-u}+\ee^{-2(t-u)}\right).
\]
Thus, with $d=|t-u|$,
\[
 F(d):=\frac12(\ee^d+\ee^{-2d})
 =\cosh d-\frac12\bigl(\ee^{-d}-\ee^{-2d}\bigr).
\]
Define
\[
 k(d):=\ee^{-|d|}-\ee^{-2|d|}.
\]
Because
\[
 \iint \cosh(t-u)\,\dd\sigma(t)\dd\sigma(u)=A[\sigma]B[\sigma],
\]
we obtain
\[
 \mathcal Q_3[\nu]
 =1-\frac{1}{2}\frac{J[\sigma]}{A[\sigma]B[\sigma]},
\]
where
\[
 J[\sigma]:=\iint_{\R^2}k(t-u)\,\dd\sigma(t)\dd\sigma(u).
\]
Therefore
\begin{equation}\label{eq:M-def}
 \beta_3=1-\frac{M}{2},
 \qquad
 M:=\sup_{\sigma\ge0}\frac{J[\sigma]}{A[\sigma]B[\sigma]}.
\end{equation}
The quotient in \eqref{eq:M-def} is invariant under multiplication of $\sigma$ by a positive scalar and under translations $t\mapsto t+c$.  Hence we may impose the two normalizations
\begin{equation}\label{eq:normalization}
 A[\sigma]=B[\sigma]=1.
\end{equation}
In this gauge,
\begin{equation}\label{eq:M-normalized}
 M=\sup\left\{J[\sigma]:\sigma\ge0,
 \int\ee^{-t}\dd\sigma=\int\ee^t\dd\sigma=1\right\}.
\end{equation}
The supremum is strictly positive.  Indeed, for any $a>0$ the two--point measure
\[
 \sigma_a:=\frac{\delta_a+\delta_{-a}}{2\cosh a}
\]
satisfies $A[\sigma_a]=B[\sigma_a]=1$ and, since $k(0)=0$ and $k(2a)>0$,
\[
 J[\sigma_a]=\frac{k(2a)}{2\cosh^2 a}>0.
\]

\begin{lemma}[Existence of a maximizer]\label{lem:existence}
The supremum in \eqref{eq:M-normalized} is attained.
\end{lemma}

\begin{proof}
Let $(\sigma_n)$ be a maximizing sequence satisfying \eqref{eq:normalization}.  The moment constraints imply the uniform tail bounds
\[
 \sigma_n((R,\infty))\le \ee^{-R},
 \qquad
 \sigma_n(( -\infty,-R))\le \ee^{-R}.
\]
Moreover, by Cauchy--Schwarz,
\[
 \sigma_n(\R)^2
 \le\left(\int\ee^t\dd\sigma_n\right)
 \left(\int\ee^{-t}\dd\sigma_n\right)=1.
\]
Hence the sequence is tight and uniformly bounded in mass.  Passing to a subsequence, $\sigma_n\rightharpoonup\sigma$ narrowly as finite positive measures.  Narrow convergence together with convergence of the total masses implies
\[
 \sigma_n\otimes\sigma_n\rightharpoonup\sigma\otimes\sigma
\]
on $\R^2$.  Since $k(t-u)$ is bounded and continuous, it follows that
\[
 J[\sigma_n]\longrightarrow J[\sigma].
\]
In particular $J[\sigma]=M>0$, so $\sigma$ is nonzero.

Lower semicontinuity of the exponential moments gives $A[\sigma]\le1$ and $B[\sigma]\le1$.  Since $\sigma$ is a nonzero positive measure, both $A[\sigma]$ and $B[\sigma]$ are strictly positive.  If $A[\sigma]B[\sigma]<1$, first multiply $\sigma$ by $(A[\sigma]B[\sigma])^{-1/2}$ and then translate it by
\[
 c=\frac12\log\frac{A[\sigma]}{B[\sigma]}.
\]
The resulting positive measure $\widetilde\sigma$ satisfies $A[\widetilde\sigma]=B[\widetilde\sigma]=1$, while amplitude scaling and translation invariance give
\[
 J[\widetilde\sigma]
 =\frac{J[\sigma]}{A[\sigma]B[\sigma]}
 >J[\sigma]=M,
\]
a contradiction to the definition of $M$.  Hence $A[\sigma]B[\sigma]=1$.  Together with $A[\sigma]\le1$ and $B[\sigma]\le1$, this forces $A[\sigma]=B[\sigma]=1$.  Thus $\sigma$ is a normalized maximizer.
\end{proof}

The variational inequality can be obtained without invoking an abstract infinite-dimensional KKT theorem.

\begin{lemma}[Obstacle condition]
Let $\sigma$ be a normalized maximizer and let $M=J[\sigma]$.  Then
\begin{equation}\label{eq:obstacle}
 2(k*\sigma)(t)\le M(\ee^t+\ee^{-t})
 \qquad\text{for all }t\in\R,
\end{equation}
and equality holds on $\supp\sigma$.
\end{lemma}

\begin{proof}
Consider $\sigma+\varepsilon\delta_t$.  Since the quotient $J/(AB)$ is invariant under subsequent amplitude and translation renormalization, maximality gives
\[
 \frac{J[\sigma+\varepsilon\delta_t]}
 {A[\sigma+\varepsilon\delta_t]B[\sigma+\varepsilon\delta_t]}
 \le M.
\]
Because $k(0)=0$, differentiation at $\varepsilon=0+$ yields \eqref{eq:obstacle}.  Define
\[
 V(t):=M(\ee^t+\ee^{-t})-2(k*\sigma)(t)\ge0.
\]
Integrating against $\sigma$ gives
\[
 \int V\,\dd\sigma=M(A+B)-2J=0.
\]
Thus $V=0$ $\sigma$-almost everywhere, and continuity of $V$ implies equality on $\supp\sigma$.
\end{proof}
